\section{Introduction}\label{sec:introduction}

Consider the scalar Allen--Cahn equation
\begin{equation}\label{eq:allen-cahn}
 \Delta u=u^3-u.
\end{equation}
The function
\[
 g(s)=\tanh(s/\sqrt2)
\]
is its increasing one-dimensional transition from $-1$ to $1$.
Thus $u(x)=g(e\cdot x-c)$ solves \eqref{eq:allen-cahn} whenever
$e$ is a unit vector; its level sets are parallel hyperplanes.
De Giorgi's conjecture asks whether every bounded entire solution
that is strictly increasing in a fixed direction has this form in
ambient dimension at most eight \cite{DeGiorgi1979}.

The conjecture connects phase transitions with the Bernstein problem
for minimal hypersurfaces.  For the quartic potential
\[
 W(s)=\frac14(1-s^2)^2,
\]
the rescaled phase-transition energies have surface tension times
perimeter as their variational limit
\cite{ModicaMortola1977,Modica1987}.
The graph condition in the minimal-surface problem is important:
nonflat minimizing cones occur in ambient dimension eight, whereas
nonaffine entire minimal graphs first occur in ambient dimension nine
\cite{Simons1968,BDG1969}.
The corresponding distinction between stable and monotone
Allen--Cahn solutions will also enter our proof.

Theorem~\ref{thm:stable} has a geometric parallel in the stable Bernstein
theorem for complete, connected, two-sided stable smooth minimal immersions
of boundaryless six-manifolds into $\R^7$: their images are entire affine
hyperplanes, with no volume-growth hypothesis
\cite[Theorem~1.1]{HongLiWang2026}.
The Bernstein theorem is not an input to our argument.

Ghoussoub and Gui established the conjecture in dimension two
\cite{GhoussoubGui1998}, and Ambrosio and Cabr\'e in dimension three
\cite{AmbrosioCabre2000}.
Savin proved the result through dimension eight under the additional
assumption
\[
 \lim_{t\to-\infty}u(x',t)=-1,\qquad
 \lim_{t\to+\infty}u(x',t)=1
 \quad\text{for every }x'\in\R^{n-1}
\]
\cite[Theorem~2.4]{Savin2009}.
These limits are pointwise in the transverse variable.
Del Pino, Kowalczyk and Wei constructed nonplanar strictly monotone
solutions in every ambient dimension at least nine, near large
dilations of nonaffine minimal graphs \cite{DPKW2011}.
The main result of this paper gives a positive resolution in
dimension eight without an assumption on the directional
limits.

\begin{theorem}[Monotone rigidity in dimension eight]
\label{thm:monotone}
Let $u\in C^2(\R^8,(-1,1))$ satisfy
\[
 \Delta u=u^3-u,\qquad \partial_8u>0\quad\text{in }\R^8.
\]
Then there are $e\in\mathbb S^7$ with $e_8>0$ and $c\in\R$ such that
\[
 u(x)=\tanh\!\left(\frac{e\cdot x-c}{\sqrt2}\right)
 \qquad\text{for every }x\in\R^8.
\]
\end{theorem}

Neither energy growth nor minimization is assumed.  Those properties,
and the well limits, will be consequences of the argument.
Both rigidity theorems concern the scalar quartic equation
\eqref{eq:allen-cahn}; the monotonicity assumption in
Theorem~\ref{thm:monotone} is strict positivity of $\partial_8u$.

\subsection{The stable equation}

For an open set $\Omega\subset\R^n$ write
\[
 \mathcal E(v;\Omega)
   =\int_\Omega\left(\frac12|\nabla v|^2+W(v)\right)\dd x.
\]
A solution $v$ is \emph{stable} if its second variation satisfies
\begin{equation}\label{eq:stability-definition}
 Q_v(\varphi):=
 \int_{\R^n}\bigl(|\nabla\varphi|^2+(3v^2-1)\varphi^2\bigr)\dd x
 \ge0
 \qquad\text{for every }\varphi\in C_c^1(\R^n).
\end{equation}
A \emph{global minimizer} minimizes $\mathcal E$ in every ball
among $H^1$ competitors with the same boundary trace.
Stability requires only the quadratic inequality, so it is a weaker
condition than minimization.

A strictly positive directional derivative solves the linearized
equation and gives stability.  Moreover, the limits along that
direction are stable solutions in one fewer variable.
This observation makes the following classification the analytic
core of the paper.

\begin{theorem}[Stable rigidity in dimension seven]
\label{thm:stable}
Let $v\in C^2(\R^7,[-1,1])$ be a stable solution of
$\Delta v=v^3-v$.  Then either $v\equiv1$, $v\equiv-1$, or
\[
 v(x)=\tanh\!\left(\frac{a\cdot x-b}{\sqrt2}\right)
 \qquad\text{for some }a\in\mathbb S^6,\ b\in\R.
\]
No energy-growth assumption is required.
\end{theorem}

This stability endpoint is sharp: Pacard and Wei constructed bounded
stable nonplanar solutions in dimension eight
\cite{PacardWei2013}.
For the smooth quartic equation, Liu, Luo, Wang, Juncheng Wei,
Yong Wei and Wu prove stable rigidity in dimension three and
monotone rigidity in dimension four \mbox{\cite{LLWWW2026}.}
The same conclusions are obtained independently by Chan, Fern\'andez-Real,
Figalli, Florit-Simon and Serra
\cite[Theorem~1.1 and Corollary~1.2]{CFFFS2026}.
Florit-Simon and Serra prove stable rigidity in dimension four
under a bounded-energy-density hypothesis \cite{FSS2025}.
Both the higher dimension and the removal of the density assumption
matter here; they require different parts of the analysis.

\subsection{From stable end profiles to the monotone theorem}

We give the complete deduction of Theorem~\ref{thm:monotone} from
Theorem~\ref{thm:stable} at the beginning of the paper.
The comparison with competitors between the directional limits is due
to Alberti, Ambrosio and Cabr\'e
\cite[Theorem~4.4]{AlbertiAmbrosioCabre2001}.
Jerison and Monneau established the minimizing-end-profile reduction
\cite[Theorem~1.3]{JerisonMonneau2004}.
The positive derivative supplies stability, and the directional
limits are stable seven-dimensional solutions.  Their classification
allows constant wells or planar transitions at this stage.
Each is a global minimizer, by the one-dimensional calibration.
Clipping a competitor between the two end profiles and sliding
translates of $u$ then shows that $u$ is itself a global minimizer.
Comparison with a constant phase inside a large ball gives
$\mathcal E(u;B_R)\le CR^7$.

The surface tension in this normalization is
\begin{equation}\label{eq:intro-tension}
 s_0=\int_{-1}^1\sqrt{2W(t)}\dd t
     =\int_\R g'(s)^2\dd s
     =\frac{2\sqrt2}{3}.
\end{equation}
Under dilations $u(R_jx)$ with $R_j\to\infty$, the minimizing
comparison identifies the entire limiting energy measure with
$s_0$ times the perimeter measure of a minimizing boundary.
This equality needs a recovery sequence patched to the original
solution; convergence of the phases alone would leave open the
possibility of extra energy.

Energy monotonicity makes the limiting boundary a cone, that is,
a boundary preserved by dilations about the origin.
Directional monotonicity fixes the sign of one component of its
normal.  In dimension eight the minimizing-boundary regularity
theory allows only a possible singularity at the vertex
\cite{Federer1970}.
The classical directed-cone argument then forces a hyperplane
\cite{Schnuerer2003}.  We include the case in which that normal
component vanishes identically: it gives translation invariance,
which is incompatible with an isolated singular vertex.
The limiting density is therefore that of one planar transition.
Wang's near-unit-density theorem
\cite[Theorem~11.1]{Wang2017} returns to the original solution,
and the one-dimensional equation determines the profile and its
orientation.

This deduction explains why the stable theorem is the right target.
It also separates the classical geometric reduction from the new
estimates needed for stable solutions.  A reader interested primarily
in the conjecture can follow the deduction first; a reader interested
in the stable classification can proceed directly from its statement
to the analytic argument in Part II.
The classical gradient bound and energy monotonicity needed by both
paths are proved in Section~\ref{sec:classical-estimates}.

\subsection{The finite-density step}

Let $\omega_{n-1}$ be the volume of the unit ball in $\R^{n-1}$.
For a solution in $\R^n$ define
\[
 M_R(y)=\frac{\mathcal E(v;B_R(y))}
              {s_0\omega_{n-1}R^{n-1}},\qquad
 M_\infty=\sup_{R>0}M_R(0).
\]
A single planar transition has limiting density one.
The first analytic step is the following result, which is also
needed locally in the argument that removes the density bound.

\begin{theorem}[Finite-density stable rigidity]\label{thm:finite}
Let $4\le n\le7$.  Suppose $v\in C^2(\R^n,[-1,1])$ solves
$\Delta v=v^3-v$ and satisfies
\[
 \int_{\R^n}\bigl(|\nabla\varphi|^2+(3v^2-1)\varphi^2\bigr)\dd x
 \ge0
 \qquad\text{for every }\varphi\in C_c^1(\R^n).
\]
If $M_\infty<\infty$, then $v\equiv1$, $v\equiv-1$, or
$v(x)=g(a\cdot x-b)$ for some $a\in\mathbb S^{n-1}$ and $b\in\R$.
\end{theorem}

For stable solutions in $\R^7$, the remaining task is to prove that
$M_\infty$ is finite. A penalized reach measures both the bending of a
zero sheet and its proximity to other sheets; the finite-density theorem
supplies local regularity at nearly this scale. Combining reach geometry
with stability bounds the fourth curvature moment and squared defect,
controlling the regions where the zero set is poorly approximated by
planes. We then compare the reach inequality with an ambient stability
test built from the same radial amplitude on the sheets. The cost of
joining neighboring amplitudes is small enough to leave a strict positive
margin, which contradicts the amount of surface area forced by unbounded
density.

\subsection{Organization and conventions}

Part I gives the complete deduction from stable end profiles to
Theorem~\ref{thm:monotone}, conditional on
Theorem~\ref{thm:stable}.
Appendix~\ref{sec:dimension-nine} explains compatibility with the
dimension-nine construction of del Pino, Kowalczyk and Wei.
After an optional guide to the analytical strategy in
Section~\ref{sec:analytic-strategy}, Part II proves the diffuse estimates,
Theorem~\ref{thm:finite}, and its local regularity consequence.
Part III develops the reach and interface estimates, joins the
ambient and surface tests, and proves fourth-moment growth.
Part IV constructs the radial ambient test and completes the
unbounded-density contradiction.
Thus the reading order places the geometric application first,
while the analytic proof of the stable theorem is contained in
Parts II--IV.  For a direct analytical reading, begin at
Section~\ref{sec:diffuse}, whose opening recalls the equation, stability
and all normalizations used there; the strategy section can be read
beforehand or consulted later.

We write $n=m+1$, so $m$ is the interface dimension, and use
$\dd S$ for induced surface measure.
The surface tension $s_0$ and the normal-average constant
$\sigma=2s_0$ are kept distinct.
All stability tests are unrestricted compactly supported tests.
Uniformity in translating centers, the order of parameter choices,
and successive limits are specified in the statements where
they are needed.
